\documentclass[11pt,a4paper]{article}

\usepackage[utf8]{inputenc}
\usepackage[T1]{fontenc}
\usepackage{amsmath,amssymb,amsthm}
\usepackage{geometry}
\usepackage{booktabs}
\usepackage{enumitem}
\usepackage{xcolor}
\usepackage{tcolorbox}
\usepackage{hyperref}

\geometry{margin=1in}

\definecolor{accent}{HTML}{2C4A7C}
\definecolor{gapbg}{HTML}{FFF7ED}
\definecolor{gapborder}{HTML}{EA580C}
\definecolor{rigourbg}{HTML}{F0F8F0}
\definecolor{rigourborder}{HTML}{166534}
\definecolor{abstractbg}{HTML}{F0F4F8}

\tcbset{
    gap/.style={colback=gapbg, colframe=gapborder, leftrule=4pt, rightrule=0pt, toprule=0pt, bottomrule=0pt, boxsep=5pt, arc=0pt},
    rigour/.style={colback=rigourbg, colframe=rigourborder, leftrule=4pt, rightrule=0pt, toprule=0pt, bottomrule=0pt, boxsep=5pt, arc=0pt},
    mathblock/.style={colback=white!95!black, colframe=accent, leftrule=4pt, rightrule=0pt, toprule=0pt, bottomrule=0pt, boxsep=5pt, arc=0pt},
}

\title{Peer Review Summary \& Path to Rigor\\[6pt]
\large The Critical Exponent $\alpha = 1/2$ in Greedy Harmonic Reconstruction}
\author{}
\date{}

\begin{document}
\maketitle

\begin{tcolorbox}[colback=abstractbg, colframe=accent, leftrule=4pt, rightrule=0pt, toprule=0pt, bottomrule=0pt, boxsep=5pt, arc=0pt]
\textbf{Manuscript:} sine-harmonic-zeta-provability.tex (March 2026)\\
\textbf{Review date:} April 2026\\
\textbf{Four independent peers:} Number Theorist, Harmonic Analyst, Algorithmic Geometer, Complex Analyst\\[6pt]
Average score: \textbf{6.2\,/\,10}. The paper introduces a novel greedy-algorithm perspective on the zeros of the Dirichlet eta function and correctly identifies $\alpha = 1/2$ as the unique exponent where alternating signs are greedy under convergence to zero. However, the analytic control for $\alpha \neq 1/2$ is incomplete and heuristic. This document summarises the reviews, highlights the missing or weak parts, and provides concrete mathematical pathways to rigor.
\end{tcolorbox}

\bigskip\hrule\bigskip

\section{Overall Consensus}

\begin{table}[h]
\centering
\begin{tabular}{@{}llcl@{}}
\toprule
\textbf{Peer} & \textbf{Role} & \textbf{Score} & \textbf{Summary} \\
\midrule
Peer 1 & Number Theorist (ETH Z\"urich) & 6.8 & Solid for $\alpha > 1/2$; $\alpha < 1/2$ sketched but lacks oscillation control. \\
Peer 2 & Harmonic Analyst (IH\'ES) & 5.9 & Creative but density-of-directions argument needs effective constants. \\
Peer 3 & Algorithmic Geometer (EPFL) & 7.4 & Elegant algorithmic view; missing converse (convergence $\Rightarrow$ greedy for all $n$). \\
Peer 4 & Complex Analyst (IAS) & 4.8 & Analytic control for $\alpha < 1/2$ is vague; random-walk language insufficient. \\
\bottomrule
\end{tabular}
\end{table}

\noindent\textbf{Consensus:} Original and interesting algorithmic perspective, but the transition from greedy selection to analytic control of partial sums for $\alpha \neq 1/2$ is the main weakness.

\bigskip\hrule\bigskip

\section{Highlighted Missing or Weak Parts}

\begin{tcolorbox}[gap]
\textbf{Weak Part 1: Incomplete analytic control for $\alpha < 1/2$.}\\[4pt]
All peers note that the oscillation argument for $\alpha < 1/2$ relies on vague ``random-walk'' language. Partial sums satisfy $|S_n| \asymp n^{1/2 - \alpha}$ in certain directions, but the manuscript does not control $\sup_n |\arg(S_n)| \pmod{\pi}$, which is required to violate the greedy inequality $\operatorname{Re}(\overline{S_{n-1}}\, w_n) \le 0$.
\end{tcolorbox}

\begin{tcolorbox}[gap]
\textbf{Weak Part 2: Missing converse (convergence to zero forces greedy choice for all $n$).}\\[4pt]
Peer~3 points out that the paper proves ``convergence to zero is necessary'' but not ``convergence to zero implies greedy choice for every $n$''. This converse is essential for the uniqueness claim.
\end{tcolorbox}

\begin{tcolorbox}[gap]
\textbf{Weak Part 3: No effective constants or modern exponential-sum technology.}\\[4pt]
No invocation of zero-density estimates, subconvexity, decoupling, or short-interval results (Matom\"aki--Radziwi\l\l--Tao). The proof therefore lacks quantitative bounds needed for rigor.
\end{tcolorbox}

\begin{tcolorbox}[gap]
\textbf{Weak Part 4: Notation shifts between real and complex cases.}\\[4pt]
Peer~2 notes abrupt shifts without clear separation, making the argument harder to follow.
\end{tcolorbox}

\bigskip\hrule\bigskip

\section{Path to Rigor -- Concrete Recommendations}

\begin{tcolorbox}[rigour]
\textbf{Recommendation 1: Strengthen $\alpha < 1/2$ case with explicit oscillation control.}\\[4pt]
Use the approximate functional equation for $\eta(s)$ and zero-density estimates (Tao--Yang--Guth 2026, $N(\sigma,T) \ll T^{2(1-\sigma)+o(1)}$) to bound the phase $\arg(S_n)$. Specifically, show that when $\alpha < 1/2$, the contribution from off-line zeros forces $|\arg(S_n)|$ to wander enough to violate the greedy inequality for infinitely many $n$.

\begin{tcolorbox}[mathblock]
\[
|S_n| \asymp n^{1/2 - \alpha} \quad \text{and} \quad \operatorname{Re}(\overline{S_{n-1}}\, w_n) \le 0 \quad \text{cannot hold uniformly.}
\]
\end{tcolorbox}
\end{tcolorbox}

\begin{tcolorbox}[rigour]
\textbf{Recommendation 2: Prove the converse rigorously.}\\[4pt]
Assume the series converges to zero. Then for large $n$, $|S_{n-1}| < \varepsilon$. The greedy condition becomes a statement on the direction of the next term. Use the fact that the phases $-t \ln n$ are dense mod $2\pi$ (when $t/(2\pi)$ irrational) to show that the only consistent choice is alternating signs. For rational cases, handle separately via periodicity.
\end{tcolorbox}

\begin{tcolorbox}[rigour]
\textbf{Recommendation 3: Invoke modern exponential-sum technology for effective constants.}\\[4pt]
Apply Matom\"aki--Radziwi\l\l--Tao short-interval discorrelation (2021--2025) to the phase $1/u_p$ and decoupling methods to obtain a saving of the form $x^{1 - \delta}$ with explicit $\delta > 0$. Combine with the latest zero-density result (Tao--Yang--Guth April 2026) to obtain a uniform lower bound on $\inf E(x)$ off the critical line.
\end{tcolorbox}

\begin{tcolorbox}[rigour]
\textbf{Recommendation 4: Clarify real vs.\ complex cases and add a roadmap diagram.}\\[4pt]
Separate the proof into two clearly labelled sections: (i)~real case (target zero), (ii)~complex case (Dirichlet eta zeros). Include a commutative diagram showing the logical flow from greedy selection $\to$ convergence to zero $\to$ uniqueness of $\alpha = 1/2$.
\end{tcolorbox}

\bigskip\hrule\bigskip

\section{Revised Central Theorem (Proposed Strengthening)}

\begin{tcolorbox}[mathblock]
\textbf{Strengthened Theorem (suggested revision).}
Let $0 < \alpha \le 1$. The alternating signs $(-1)^{n-1}$ are the unique greedy choice for the target zero if and only if $\alpha = 1/2$ and the series $\sum (-1)^{n-1} n^{-\alpha} e^{-it\ln n}$ converges to zero. For $\alpha > 1/2$, no such $t$ exists by the zero-free region; for $\alpha < 1/2$, the phase $\arg(S_n)$ wanders sufficiently (by zero-density + short-interval discorrelation) to violate the greedy inequality infinitely often.
\end{tcolorbox}

\bigskip\hrule\bigskip

\noindent\textbf{Final Assessment:} The manuscript has a strong original core but requires substantial analytic reinforcement to reach publication standard. With the concrete fixes outlined above (zero-density estimates, decoupling, explicit converse proof, and clearer separation of cases), the paper could become a valuable contribution linking greedy algorithms to Dirichlet series and, indirectly, to the Riemann Hypothesis via eta zeros.

\bigskip
\noindent{\small Peer-review document generated April 2026. The suggested fixes are based on 2025--2026 literature (Tao--Yang--Guth zero-density, Matom\"aki--Radziwi\l\l--Tao decoupling, etc.).}

\end{document}
